% ============================================================
%  appendix_repairs.tex : Biology-faithful overlay rules
%
%  Transcribed verbatim from the model source
%  setup_a/bbcn/honest_bbcn.py (step_honest, patient_clamp_off),
%  covering R1 (p53-MDM2-ATM), R2 (AKT1-FOXO3a-PHLPP + flag U),
%  and R3 (the explicit caspase-cascade commitment mechanism).
%
%  Hand-maintained against the code; NOT the auto-generated
%  appendix_rules.tex (which lists the BARE 135-node network).
%
%  Durability engine: the reported durable-fixed-point figures are
%  certified on the free caspase cascade (step_honest(cascade=True)),
%  no node forced. A fully cascade-native controller (kernel search
%  also on the cascade) is deferred to post-resubmission work and
%  recovers additional fixed points (~20%); the conservative figure
%  below is the subset durable under BOTH engines.
% ============================================================

\section{The biology-faithful overlay (repaired rules)}
\label{app:repairs}

Appendix~\ref{app:rules} lists the baseline 135-node network verbatim from
\texttt{pathways.py}. The biology-faithful model applies a runtime overlay that
leaves \texttt{pathways.py} untouched, replaces the update rules of four baseline
nodes, and adds six nodes (PHLPP, then PUMA, NOXA, SMAC, CASP6, CASP7), giving a
\textbf{141-node} signalling model. It also introduces a per-patient uncoupling flag
$U$. The rules below are transcribed from \texttt{setup\_a/bbcn/honest\_bbcn.py}
(\texttt{step\_honest}, \texttt{patient\_clamp\_off}); where a rule replaces a
baseline rule, the baseline form of Appendix~\ref{app:rules} no longer applies.

\paragraph{R1. p53, MDM2, and ATM damage sensor.}\mbox{}\\
The baseline arm ($\mathrm{MDM2}=\mathrm{TP53}\lor\mathrm{MDM2}$, self-locking, and
$\mathrm{TP53}=\lnot\mathrm{MDM2}$, no damage input) is replaced by a sensor in which
DNA damage stabilises p53. The damage signal unions the endogenous sensors (ATM/ATR)
with the genotoxic treatment input (default $0$):
\begin{flushleft}\ttfamily\small
dmg = ATM OR ATR OR genotoxic\\
MDM2 = (TP53 OR AKT1) AND NOT (CDKN2A AND (E2F1 OR dmg))\\
TP53 = (NOT MDM2) OR dmg\\
\end{flushleft}

\paragraph{R2. AKT1, FOXO3a, and PHLPP loop and the uncoupling flag $U$.}\mbox{}\\
The baseline AKT and FOXO arm ($\mathrm{AKT1}=\mathrm{MTOR}\land\mathrm{PDPK1}\land
\mathrm{PIK3CA}\land\lnot\mathrm{PTEN}$; $\mathrm{FOXO3}=\lnot\mathrm{AKT1}$) has no
closing feedback. The repair adds PHLPP (the phosphatase that closes the loop) and a
per-patient flag $U$ (\texttt{CLAMP\_OFF}), computed once at initialisation from the
binarised state and held constant: $U=1$ when a strict majority of the AKT-activity
nodes and both FOXO-nuclear nodes are ON.
\begin{flushleft}\ttfamily\small
U = 1 if ( |\{AKT1,MTOR,RPS6KB1,EIF4EBP1\} ON| $>$ 2  AND  |\{FOXO3,FOXO1\} ON| $=$ 2 ) else 0\\
PHLPP = FOXO3 AND NOT U\\
AKT1 = ((MTOR AND PDPK1 AND PIK3CA AND NOT PTEN) AND NOT PHLPP) OR (U AND AKT1)\\
FOXO3 = (NOT AKT1) OR U\\
\end{flushleft}
\noindent The term \texttt{(U AND AKT1)} is the self-sustaining uncoupled latch: under
$U=1$, AKT1 stays high once high. This is an \emph{encoding} of the uncoupled-resistance
configuration, not an emergent finding; the model then tests which patients survival-axis
inhibition can flip (the coupled fraction), which is the reported result. PHLPP is
dynamically inert outside the switch, so the strongly-connected-component and feasibility
analyses of Section~\ref{sec:obstruction} are identical on the 135- and 136-node
models.

\paragraph{R3. Commitment by the intrinsic caspase cascade.}\mbox{}\\
Commitment is the explicit p53$\to$BH3-only$\to$executioner cascade, not an abstract
brake-zeroing flag. p53 induces PUMA and NOXA, which repress the anti-apoptotic
guardians; BAX/BAK release SMAC, which neutralises XIAP; and an executioner loop makes
death irreversible. The five nodes PUMA, NOXA, SMAC, CASP6, CASP7 are added here (CASP8,
CASP3, and the guardians already exist in the baseline), taking the model to 141 nodes.
\begin{flushleft}\ttfamily\small
PUMA = TP53\\
NOXA = TP53\\
MCL1 = MCL1 AND NOT NOXA\\
BCL2 = BCL2 AND NOT PUMA\\
BCL2L1 = BCL2L1 AND NOT PUMA\\
SMAC = BAX OR BAK1\\
XIAP = XIAP AND NOT SMAC\\
CASP7 = (CASP8 OR CASP9) AND NOT XIAP\\
CASP6 = CASP3 OR CASP7\\
CASP8 = CASP8 OR CASP6\\
CASP3 = (CASP8 OR CASP9) AND NOT XIAP\\
\end{flushleft}
\noindent The nodes resolve within the step in the order shown, so the final line closes
the loop $\mathrm{CASP6}\to\mathrm{CASP8}\to\mathrm{CASP3}\to\mathrm{CASP6}$: once an
executioner fires it self-sustains, which is the molecular point of no return. No node is
forced; the guardians fall because PUMA/NOXA/SMAC repress them through these rules.

\paragraph{Durability readout.}\mbox{}\\
A patient is scored durable when the free dynamics of the 141-node cascade settle to a
genuine fixed point with $\mathrm{CASP3}=1$ (one synchronous update leaves the state
unchanged; no node held). Across the three cohorts this holds for
\textbf{17.1\% (TCGA), 14.6\% (METABRIC), 16.4\% (I-SPY2)} of patients, where the
baseline 135-node network supports no durable apoptotic fixed point at all (0\%). Every
reported fixed point is certified under the free cascade; a fully cascade-native kernel
search recovers a further set (raising the figure to $\sim$20\%) and is future work.
